\chapter{Quando la paura salvava la vita}

\begin{quotation}
``La paura è il prezzo che paghiamo per l'anticipazione.''\\
---\index[main]{LeDoux, Joseph}Joseph LeDoux
\end{quotation}

\section{L'eredità del terrore}

Centomila anni fa, nella savana africana. Il sole tramonta. Un nostro antenato si allontana dal gruppo per raccogliere radici. Un fruscio nell'erba alta. Potrebbe essere il vento. Potrebbe essere una preda. Potrebbe essere un leopardo. Ha una frazione di secondo per decidere.

Chi si fermava ad analizzare razionalmente la situazione -- valutare la direzione del vento, l'altezza dell'erba, la probabilità statistica di predatori nella zona -- non ha lasciato discendenti. Chi saltava e correva al minimo rumore sospetto ha tramandato i suoi geni. E con essi, un sistema di allarme calibrato su un principio brutale: meglio mille spaventi inutili che una sola morte evitabile\footnote{Marks, I. M., \& Nesse, R. M. (1994). ``Fear and fitness: An evolutionary analysis of anxiety disorders''. \textit{Ethology and Sociobiology}, 15(5-6), 247-261.}.

Oggi, quel sistema di allarme ancestrale vive ancora in voi. Solo che adesso non reagisce ai fruscii nella savana ma alle vibrazioni del telefono, ai toni di voce del capo, ai titoli dei giornali online. Il corpo si prepara a combattere o fuggire da minacce che non richiedono né combattimento né fuga: un'email, una notifica, un messaggio letto ma senza risposta\footnote{Sapolsky, R. M. (2004). \textit{Why Zebras Don't Get Ulcers}. New York: Times Books.}. 

Questo è il paradosso della paura moderna: siamo \index[main]{cacciatori-raccoglitori}cacciatori-raccoglitori biologici in un mondo di informazioni digitali, e il nostro \index[main]{sistema di allarme}sistema di allarme non ha ancora imparato la differenza tra un pixel e un predatore.


\section{Il genitore iperprotettivo che vive in noi}

\index[main]{Nesse, Randy}Randy Nesse ha capito che tutta questa logica ancestrale segue un principio elegante nella sua semplicità: quello che lui chiama il \index[main]{smoke detector principle}smoke detector principle''\footnote{Nesse, R. M. (2005). Natural selection and the regulation of defenses: A signal detection analysis of the smoke detector principle''. \textit{Evolution and Human Behavior}, 26(1), 88-105.}.

Per comprenderlo, basta osservare il comportamento di un genitore con il primo figlio.
Ogni colpo di tosse diventa un possibile polmonite. Ogni caduta, una potenziale commozione cerebrale. Ogni febbriciattola, un'emergenza pediatrica. Gli altri genitori sorridono indulgenti -- sanno che al terzo figlio controllerai a malapena se respira ancora. Ma il meccanismo è perfetto nella sua apparente irrazionalità: meglio dieci corse inutili al pronto soccorso che perdere l'unica volta in cui c'è davvero qualcosa che non va.
Il vostro sistema immunitario funziona con la stessa logica spietata. Reagisce in modo eccessivo a pollini innocui, vi fa starnutire per una settimana per quello che potrebbe essere solo polvere, scatena l'inferno delle allergie per sostanze che non vi farebbero nulla. Ma quella stessa ipersensibilità vi protegge da virus e batteri che potrebbero uccidervi. Il sistema preferisce bombardare tutto piuttosto che rischiare di lasciar passare il nemico vero.
Anche il vostro antivirus funziona così: vi blocca siti perfettamente innocui, vi impedisce di scaricare file che avreste voluto, vi tempesta di avvisi per ogni cosa. Esasperante? Sì. Ma preferibile all'alternativa di un sistema permissivo che lascia entrare il malware che vi distrugge l'hard disk.
L'evoluzione ha applicato questa stessa filosofia ai nostri sistemi di sopravvivenza emotiva. Solo che adesso viviamo in un mondo dove il "genitore iperprotettivo" che abbiamo in testa scambia una email del capo per un attacco di tigre dai denti a sciabola. E a differenza dei veri genitori, che col tempo imparano a rilassarsi, il nostro sistema di allarme ancestrale non ha mai imparato che alcune minacce erano solo per finta\footnote{Haselton, M. G., & Nettle, D. (2006). ``The paranoid optimist: An integrative evolutionary model of cognitive biases''. \textit{Personality and Social Psychology Review}, 10(1), 47-66.}.
Questo ci porta a una domanda inquietante: cosa succede quando un sistema progettato per reagire a pericoli fisici immediati deve gestire minacce psicologiche che non finiscono mai? La risposta è scritta nelle prescrizioni di ansiolitici del mondo occidentale...

\section{La fabbrica delle emozioni}

Ma l'\index[main]{amigdala}amigdala e il suo sistema di allarme sono solo la punta dell'iceberg. Il nostro cervello è una vera e propria \index[main]{fabbrica delle emozioni}fabbrica di emozioni, con reparti specializzati, catene di montaggio complesse, e un sistema di controllo qualità che spesso va in tilt. È un'operazione industriale su scala neuronale, con milioni di cellule che lavorano 24 ore su 24 per produrre quella miscela chimica che chiamiamo ``sentimenti''.

Se il cervello fosse un'azienda, le emozioni sarebbero il reparto marketing: influenzano ogni decisione, spesso in modi di cui non siamo consapevoli, e hanno la tendenza a promettere cose che il reparto ``razionalità'' poi non riesce a mantenere.

Il problema è che questa fabbrica è stata progettata milioni di anni fa per produrre un catalogo di emozioni molto limitato -- paura, rabbia, gioia, tristezza, disgusto, sorpresa -- ma oggi deve gestire ordini molto più complessi: ansia da prestazione, invidia sui social media, \index[main]{FOMO}FOMO, burnout professionale, crisi esistenziali\footnote{Barrett, L. F. (2017). \textit{How Emotions Are Made: The Secret Life of the Brain}. Boston: Houghton Mifflin Harcourt.}.

È come chiedere a una fabbrica di candele del 1800 di produrre smartphone: tecnicamente è sempre una questione di assemblare componenti, ma la complessità è aumentata in modo esponenziale mentre i macchinari sono rimasti gli stessi.

\section{Il catalogo ancestrale delle paure}

\index[main]{Seligman, Martin}Martin Seligman ha scoperto che esistono ``\index[main]{paure preparate}paure preparate'' -- quelle per cui siamo biologicamente predisposti -- e ``\index[main]{paure apprese}paure apprese''\footnote{Seligman, M. E. (1971). ``Phobias and preparedness''. \textit{Behavior Therapy}, 2(3), 307-320.}. 

Il nostro kit include serpenti, ragni, altezze, spazi chiusi, sangue, estranei, isolamento sociale\footnote{Öhman, A., \& Mineka, S. (2003). ``The malicious serpent: Snakes as a prototypical stimulus for an evolved module of fear''. \textit{Current Directions in Psychological Science}, 12(1), 5-9.}. Notate qualcosa? Niente automobili nella lista, anche se uccidono migliaia di persone all'anno. Niente prese elettriche, niente armi da fuoco. Il nostro catalogo è fermo al Pleistocene.

Il ``\index[cognitive]{one-trial fear conditioning}one-trial fear conditioning'' significa che basta una singola esperienza negativa per creare una fobia permanente\footnote{Öhman, A., \& Mineka, S. (2001). ``Fears, phobias, and preparedness: Toward an evolved module of fear and fear learning''. \textit{Psychological Review}, 108(3), 483-522.}. Un morso di cane nell'infanzia, paura dei cani per sempre. Il cervello ha due standard: una prova per certificare il pericolo, infinite prove insufficienti per certificare la sicurezza.

\section{Il sistema della ricompensa hackerato}

Ma la paura è solo uno dei prodotti della fabbrica emotiva. Prendiamo il \index[main]{sistema dopaminergico}sistema della ricompensa -- quello che i neuroscienziati chiamano ``circuito dopaminergico''. È il reparto della fabbrica responsabile per la motivazione, il piacere, e la dipendenza. Il suo slogan aziendale potrebbe essere: ``Se ti fa sentire bene, deve essere importante per la sopravvivenza''\footnote{Schultz, W. (2015). ``Neuronal reward and decision signals: From theories to data''. \textit{Physiological Reviews}, 95(3), 853-951.}.

Nella versione originale, questo sistema funzionava benissimo. Cibo? Rilascio di dopamina. Sesso? Dopamina a fiumi. Successo sociale? Pioggia di dopamina. Era un sistema semplice ed efficace: quello che ti faceva sentire bene era generalmente quello che aumentava le tue possibilità di sopravvivenza e riproduzione.

Ma poi sono arrivati droghe, gioco d'azzardo, social media, cibo ultraprocessato, shopping online, pornografia, videogiochi. Tutte cose che attivano il sistema della ricompensa ma non hanno nulla a che fare con la sopravvivenza evolutiva\footnote{Volkow, N. D., Wang, G. J., Fowler, J. S., \& Telang, F. (2008). ``Overlapping neuronal circuits in addiction and obesity''. \textit{Philosophical Transactions of the Royal Society B}, 363(1507), 3191-3200.}. È come se qualcuno avesse hackerato il sistema operativo della fabbrica.

Il risultato è un sistema in perpetuo corto circuito. Il cervello continua a dire: ``Questo TikTok deve essere importante per la tua sopravvivenza, altrimenti perché ti starebbe dando tanta soddisfazione?'' E voi continuate a scrollare, mentre la parte razionale del cervello urla disperatamente: ``Ma che cosa stiamo facendo?!''

\section{L'alchimia delle emozioni complesse}

L'aspetto più affascinante della fabbrica emotiva è come produce emozioni che sembrano semplici ma sono in realtà incredibilmente sofisticate. Prendiamo quella che chiamiamo ``\index[main]{gelosia}gelosia''. Sembra un'emozione singola, unitaria. In realtà, è un cocktail complesso di almeno sei componenti emotive diverse.

C'è la paura dell'abbandono (prodotta dall'amigdala), la rabbia per la perdita di controllo (corteccia prefrontale ventromediale), l'ansia per l'incertezza (lobo dell'insula), la tristezza per la perdita percepita (sistema serotoninergico), l'invidia per quello che ha l'altro (corteccia cingolata anteriore), e perfino una componente di disgusto morale (corteccia orbitofrontale)\footnote{Harmon-Jones, E., Peterson, C. K., \& Harris, C. R. (2009). ``Jealousy: Novel methods and neural correlates''. \textit{Emotion}, 9(1), 113-117.}.

Tutti questi ``ingredienti'' vengono miscelati insieme in proporzioni diverse a seconda della situazione, della persona, e perfino dell'ora del giorno. E tutto questo succede in millisecondi, senza alcuna supervisione cosciente. È come se la fabbrica emotiva avesse un sistema di automazione così sofisticato che può produrre emozioni complesse mentre voi state pensando a tutt'altro.

\section{I difetti di produzione: quando la fabbrica sbaglia}

Come tutte le fabbriche, anche quella emotiva ha i suoi difetti di produzione. E alcuni di questi difetti sono diventati i nostri bias cognitivi più problematici.

Prendiamo quello che gli psicologi chiamano ``\index[cognitive]{affect heuristic}affect heuristic'' -- la tendenza a giudicare situazioni complesse in base a come ci fanno sentire piuttosto che in base a un'analisi razionale dei fatti\footnote{Slovic, P., Finucane, M. L., Peters, E., \& MacGregor, D. G. (2007). ``The affect heuristic''. \textit{European Journal of Operational Research}, 177(3), 1333-1352.}. Un investimento vi ``sembra'' giusto? Deve essere un buon investimento. Una persona vi dà una ``brutta sensazione''? Deve essere una persona cattiva.

Il problema è che le nostre emozioni sono calibrate per problemi di 50.000 anni fa, non per le complessità del mondo moderno. È per questo che tendiamo a essere più generosi con le persone attraenti (il cervello emotivo interpreta l'attrattiva come segno di buoni geni), a fidarci di più di chi ci assomiglia (similarità = appartenenza al gruppo = sicurezza), e a essere più sospettosi verso chi parla con accento straniero (diversità linguistica = potenziale minaccia esterna).

\section{La produzione emotiva su commissione}

C'è un difetto di produzione particolarmente interessante che potremmo chiamare ``\index[main]{emotional manufacturing}emotional manufacturing'' -- la produzione di emozioni su commissione.

Avete mai notato come, quando siete già di cattivo umore, tutto sembra andare storto? Non è che il mondo sia diventato improvvisamente più ostile; è che la vostra fabbrica emotiva, già settata su ``modalità negativa'', inizia a produrre interpretazioni negative di eventi neutri\footnote{Bower, G. H. (1981). ``Mood and memory''. \textit{American Psychologist}, 36(2), 129-148.}.

Un collega che non vi saluta? Deve essere arrabbiato con voi (invece che semplicemente distratto). Il commesso che è un po' brusco? Ce l'ha personalmente con voi (invece che avere semplicemente una giornata no). Il vostro partner che risponde monosillabi? Vi sta punendo per qualcosa (invece che essere semplicemente stanco).

È come se la fabbrica emotiva, una volta avviata in una certa direzione, iniziasse a produrre prove a sostegno di quella direzione. È il bias di conferma applicato alle emozioni: non solo cerchiamo informazioni che confermano quello che pensiamo, ma produciamo anche emozioni che confermano quello che sentiamo.

\section{Il sequestro emotivo}

Il difetto più problematico è quello che potremmo chiamare ``\index[main]{emotional hijacking}emotional hijacking'' -- il sequestro emotivo. È quello che succede quando l'amigdala decide che c'è un'emergenza e prende il controllo di tutta l'operazione\footnote{Goleman, D. (1995). \textit{Emotional Intelligence}. New York: Bantam Books.}.

In questi momenti, la fabbrica emotiva smette di produrre la normale gamma di sensazioni bilanciate e inizia a sfornare massicce quantità di una singola emozione -- di solito rabbia o paura. Siete ``troppo arrabbiati per pensare'', ``troppo spaventati per ragionare''. La parte logica del cervello è letteralmente offline.

Questo meccanismo aveva senso quando le emergenze erano leoni che attaccavano. Ma oggi le nostre ``emergenze'' sono più spesso email del capo o commenti sui social. Situazioni che richiederebbero una risposta misurata, ma che invece attivano il protocollo di emergenza ancestrale.

\section{Il paradosso della sicurezza ansiosa}

Viviamo oggettivamente nell'epoca più sicura della storia umana. \index[main]{Pinker, Steven}Steven Pinker ha documentato il declino drastico di quasi ogni forma di violenza\footnote{Pinker, S. (2011). \textit{The Better Angels of Our Nature: Why Violence Has Declined}. New York: Viking.}. Eppure i disturbi d'ansia colpiscono il 31\% degli americani nel corso della vita\footnote{Kessler, R. C., et al. (2005). ``Lifetime prevalence and age-of-onset distributions of DSM-IV disorders''. \textit{Archives of General Psychiatry}, 62(6), 593-602.}. L'uso di ansiolitici è aumentato del 400\% dal 1988\footnote{Bachhuber, M. A., et al. (2016). ``Increasing benzodiazepine prescriptions and overdose mortality''. \textit{American Journal of Public Health}, 106(4), 686-688.}.

\index[main]{Ropeik, David}David Ropeik suggerisce che questa non è una contraddizione ma una conseguenza diretta\footnote{Ropeik, D. (2010). \textit{How Risky Is It, Really?}. New York: McGraw-Hill.}. In assenza di tigri reali, il sistema di allarme abbassa la soglia e inizia a reagire a tigri sempre più metaforiche.

Ma c'è un elemento che rende tutto peggiore: la persistenza. Nel mondo ancestrale, le minacce apparivano e scomparivano. Oggi, le ``minacce'' sono croniche\footnote{McEwen, B. S. (2017). ``Neurobiological and systemic effects of chronic stress''. \textit{Chronic Stress}, 1, 2470547017692328.}. Il feed di notizie negative è infinito. I social media forniscono allarmi 24/7. L'email ansiogena resta nella inbox come un predatore che non se ne va mai.

\section{La velocità del terrore, la lentezza della ragione}

Quando il sistema rileva una minaccia, l'attivazione avviene in 12-20 millisecondi\footnote{LeDoux, J. (1996). \textit{The Emotional Brain}. New York: Simon \& Schuster.}. La corteccia prefrontale -- lo scienziato all'attico che potrebbe dire ``calma, è solo un'email'' -- impiega almeno 250 millisecondi per entrare in azione\footnote{Bar, M., et al. (2006). ``Top-down facilitation of visual recognition''. \textit{PNAS}, 103(2), 449-454.}.

\begin{figure}[h!]
    \centering
    \includegraphics[width=0.8\textwidth]{cap3_amig.png}
    \caption{La gara persa in partenza: il gap temporale tra emozione e ragione}
    \label{fig:funzionamento_amigdala}
\end{figure}

Per un quarto di secondo siete letteralmente in balia della reazione primitiva. È come se due piloti guidassero la stessa macchina, ma uno avesse il controllo esclusivo per i primi metri di ogni curva.

\section{Supervisori di una fabbrica antica}

Alla fine, dobbiamo accettare una verità scomoda: le emozioni non sono qualcosa che ``ci capita''. Sono qualcosa che il nostro cervello produce attivamente, costantemente, in risposta a stimoli interni ed esterni. Siamo tutti dei piccoli produttori industriali di emozioni, con fabbriche personali che lavorano senza sosta.

E come tutte le fabbriche, la qualità del prodotto dipende dalla qualità delle materie prime che usiamo. Se alimentiamo la nostra fabbrica emotiva con stress cronico, mancanza di sonno, cibo spazzatura, relazioni tossiche, e sovrastimolazione mediatica, non dovremmo stupirci se il prodotto finale è di qualità scadente.

Tecniche come la meditazione, la terapia cognitivo-comportamentale, l'esercizio fisico regolare possono essere viste come forme di ``manutenzione industriale'' per la nostra fabbrica emotiva\footnote{Davidson, R. J., \& McEwen, B. S. (2012). ``Social influences on neuroplasticity: Stress and interventions to promote well-being''. \textit{Nature Neuroscience}, 15(5), 689-695.}. Non eliminano i bias o i difetti di produzione -- quelli sono parte dell'hardware -- ma possono aiutare a ottimizzare i processi.

Il trucco è smettere di aspettarsi che la fabbrica emotiva produca solo emozioni ``razionali''. È una fabbrica antica che usa macchinari progettati per un mondo diverso. Il massimo che possiamo fare è imparare a essere supervisori consapevoli della nostra produzione emotiva.

Per quanto ci piaccia pensare di essere esseri razionali che occasionalmente provano emozioni, la realtà è più vicina al contrario: siamo esseri emotivi che occasionalmente riescono a essere razionali.

\bigskip

Nel prossimo capitolo esploreremo proprio questa convivenza forzata. Scopriremo come quello che \index[main]{Kahneman, Daniel}Daniel Kahneman ha chiamato Sistema 1 e Sistema 2 non sono solo due modi di pensare, ma due veri e propri personaggi che abitano la nostra mente, ciascuno con le proprie regole, i propri obiettivi, e la propria visione del mondo\footnote{Kahneman, D. (2011). \textit{Thinking, Fast and Slow}. New York: Farrar, Straus and Giroux.}.

Vedremo come il Sistema 1 -- quello che gestisce questa fabbrica emotiva -- prende il controllo della maggior parte delle nostre decisioni quotidiane, mentre il Sistema 2 -- quello che dovrebbe supervisionare -- preferisce sonnecchiare sul divano mentale. E scopriremo perché, paradossalmente, più il mondo diventa complesso e richiede analisi razionale, più tendiamo ad affidarci al pilota automatico emotivo.